% ===== Cómo funcionaría (B16, adaptada) =====
\begin{frame}{Cómo funcionaría}
  \vspace{0.5em}
  \begin{itemize}\setlength{\itemsep}{0.35em}
    \item Quien tiene empleo formal sigue pagando sus \clave{cuotas obrero-patronales}, como hoy. Lo que cambia es que debe ver con claridad qué paga y qué recibe.
    \item Quien no cotiza se inscribe en los seguros médicos del IMSS y paga una \clave{prima según su ingreso}. El gobierno completa la diferencia.
    \item Sólo a la población más pobre el gobierno le paga la prima completa, como un impuesto negativo. Si la pagara para todos, se perdería el incentivo a cotizar.
  \end{itemize}
  \vspace{-0.2em}
  \banda{Todos dentro del mismo sistema de aseguramiento, sin esquemas paralelos.}
  \vspace{-0.9em}
  \banda{\small La propuesta no sube las cuotas obrero-patronales. Con el cambio en Infonavit, la tasa legal sobre la nómina bajaría 3 puntos.}
\end{frame}
